%============================================================ 교전 시나리오 표 (국방발표 판)
%  저널판은 tables/tab_scenario.tex(자동 생성)를 쓰고 길이를 관측 픽셀 한 변 $\ell$ 단위로
%  적는다. 발표판은 같은 설정을 m·m/s 로 환산해 싣는다 --- 환산 기준은 기준 인스턴스
%  $\ell = 252\unit{m}$, 결정 주기 25\,s 이며, 원값은 학습 체크포인트의 설정이다.
%    교전 해역 한 변 50 \ell = 12.6 km,  적 0.89 \ell/주기 = 9.0 m/s,  방어정 0.60 \ell/주기 = 6.0 m/s
%  ★ 15쪽판은 물리 단위(m·m/s·s)를 싣지 않는다 --- 길이는 관측 픽셀 한 변 ell 단위,
%    속력은 비율만 둔다. 결정 주기·픽셀 크기의 절대값은 위 환산 기준에만 남긴다.
%  표 밑 주석은 싣지 않는다(발표판 방침).
\begin{table}[htbp]
\centering
\caption{교전 시나리오와 관측 설정}
\label{tab:scenario}
\footnotesize
\resizebox{\ifdim\width>\linewidth\linewidth\else\width\fi}{!}{%
\begin{tabular}{lr}
\toprule
항목 & 값 \\
\midrule
적 USV 수 $M$ & 10 \\
방어정 수 $P$ & 3 \\
적 무리 최대 수 $C$ & 4 \\
교전 해역 한 변 & $50\,\ell$ \\
속력비 $v_e/v_a$ (적 USV / 방어정) & 1.5 \\
관측 격자 $n_{\mathrm{g}} \times n_{\mathrm{g}}$ (모선 중심) & $50 \times 50$ \\
방어정당 그물 (학습 / 평가) & 1 / 3 \\
정책 파라미터 수 & 244{,}389 \\
\bottomrule
\end{tabular}
}
\end{table}
